\documentclass[12pt]{article}
\usepackage[letterpaper,top=0.75in,bottom=0.75in,left=0.78in,right=1.55in,marginparwidth=1.15in,marginparsep=0.18in]{geometry}
\usepackage{fontspec}
\setmainfont{Latin Modern Roman}
\usepackage{microtype}
\usepackage{fancyhdr}
\usepackage{titlesec}
\usepackage{enumitem}
\usepackage{xcolor}
\usepackage[hidelinks]{hyperref}
\setlength{\parindent}{1.05em}
\setlength{\parskip}{0.32em}
\setlength{\headheight}{15pt}
\titleformat{\section}{\large\bfseries\scshape}{}{0pt}{}
\titleformat{\subsection}{\normalsize\bfseries}{}{0pt}{}
\pagestyle{fancy}
\fancyhf{}
\fancyhead[L]{Whately Logic: Julius Caesar Lit-Anchor}
\fancyhead[R]{Lesson 5}
\fancyfoot[C]{\thepage}
\renewcommand{\headrulewidth}{0.4pt}

\newcommand{\focusbox}[1]{\par\vspace{0.35em}\noindent\begingroup\setlength{\fboxsep}{6pt}\colorbox{gray!12}{\begin{minipage}{\dimexpr\linewidth-2\fboxsep\relax}#1\end{minipage}}\endgroup\par\vspace{0.45em}}
\newcommand{\studentline}{\par\vspace{0.68em}\hrule\vspace{0.72em}}
\newcommand{\shortline}{\par\vspace{0.42em}\hrule\vspace{0.55em}}
\newcommand{\gloss}[1]{\marginpar{\scriptsize\raggedright\setlength{\baselineskip}{7.8pt}\color{gray!70!black}#1}}
\newcommand{\speaker}[1]{\par\vspace{0.35em}\noindent\textsc{#1}\par\vspace{-0.15em}}

\begin{document}
\begin{center}
{\LARGE\bfseries Julius Caesar Lit-Anchor}\par\vspace{0.18em}
{\large\scshape Brutus, Antony, and the Shape of Proof}\par\vspace{0.35em}
{Lesson 5: Mood, Figure, and Formal Validity}\par\vspace{0.55em}
\rule{0.82\textwidth}{0.4pt}
\end{center}

\section*{Purpose}
A lit-anchor connects the week's logic lesson to a recurring literary text. In Lesson 5, Whately asks us to notice that arguments have shape. The topic may be friendship, honor, ambition, slavery, or Rome, but the form still matters. A conclusion may sound noble and still fail formally. A conclusion may be emotionally painful and still follow from its premises.

\textit{Julius Caesar} gives us a public test case. Brutus explains why Caesar had to die. Antony answers him before the same crowd. Both men speak about honor and ambition, but their speeches depend on different forms of reasoning. This handout asks whether the shape of each argument actually carries its conclusion.

\focusbox{\textbf{Whately's Lesson 5 habit:} Identify the conclusion, name the A/E/I/O mood, locate the middle term, and ask whether the figure of the argument lets the conclusion follow. Do the form test before the fact test.}

\section*{Central Question}
Does the shape of Brutus's public argument prove that Caesar deserved death, or does Antony expose a weakness in its form or premises?

\section*{Guided Text: Brutus in Act 3, Scene 2}
Brutus has helped kill Caesar. He now speaks to the Roman citizens. Read the original text below. The side notes are light glosses only; they do not replace the text.

\speaker{Brutus}
\begin{verse}
Romans, countrymen, and lovers! hear me for my\\
cause, and be silent, that you may hear: believe me\\
for mine honour, and have respect to mine honour, that\\
you may believe: censure me in your wisdom, and\\
awake your senses, that you may the better judge.\gloss{Brutus asks the crowd to judge him by honor and wisdom.}\\
If there be any in this assembly, any dear friend of\\
Caesar's, to him I say, that Brutus' love to Caesar\\
was no less than his. If then that friend demand\\
why Brutus rose against Caesar, this is my answer:\\
--Not that I loved Caesar less, but that I loved\\
Rome more. Had you rather Caesar were living and\\
die all slaves, than that Caesar were dead, to live\\
all free men?\gloss{Brutus frames the choice as Caesar alive with slavery or Caesar dead with freedom.} As Caesar loved me, I weep for him;\\
as he was fortunate, I rejoice at it; as he was\\
valiant, I honour him: but, as he was ambitious, I\\
slew him. There is tears for his love; joy for his\\
fortune; honour for his valour; and death for his\\
ambition.\gloss{This line gives the argument's key middle term: ambitious.} Who is here so base that would be a\\
bondman? If any, speak; for him have I offended.\\
Who is here so rude that would not be a Roman? If\\
any, speak; for him have I offended. Who is here so\\
vile that will not love his country? If any, speak;\\
for him have I offended. I pause for a reply.
\end{verse}

\section*{Guided Analysis}
Brutus's argument can be expressed as a simple form. The form may be valid even if the premises still need to be tested.

\subsection*{A possible formal map}
\begin{itemize}[leftmargin=*]
\item \textbf{Major premise:} All men whose ambition would enslave Rome deserve death for Rome's freedom.
\item \textbf{Minor premise:} Caesar was a man whose ambition would enslave Rome.
\item \textbf{Conclusion:} Therefore, Caesar deserved death for Rome's freedom.
\item \textbf{Mood:} A-A-A.
\item \textbf{Figure pattern:} All M are P; All S are M; therefore All S are P.
\item \textbf{Middle term:} men whose ambition would enslave Rome.
\end{itemize}

\subsection*{The form test}
This is a valid beginner form. If every M belongs inside P, and if S belongs inside M, then S belongs inside P. In that narrow sense, the conclusion follows from the premises.

\subsection*{The fact test}
The form test does not settle the whole matter. Brutus must still prove the minor premise: Caesar's ambition would enslave Rome. The crowd may accept the word \textit{ambitious} because Brutus is honorable, because the fear of slavery is powerful, or because the speech gives them only two choices. Whately's form test shows where the argument bears weight.

\focusbox{\textbf{Lesson 5 takeaway:} Brutus's public argument has a recognizable shape. Its formal strength does not automatically prove its factual premise. A valid form can still depend on a contested claim.}

\section*{Discussion}
\begin{enumerate}[leftmargin=*]
\item What conclusion does Brutus want the citizens to accept?
\item What is the middle term in the formal map above?
\item Does Brutus's form appear valid or invalid? Explain using S, P, and M.
\item Which premise needs the most proof?
\item How does Brutus's question about slavery make the conclusion feel morally urgent?
\item Why is it dangerous to confuse a valid form with a fully proved argument?
\end{enumerate}

\section*{Independent Text: Antony in Act 3, Scene 2}
Antony answers Brutus without directly calling him a liar. He repeats Brutus's claim and then offers examples that put pressure on the word \textit{ambitious}. Read the original text. Then decide what formal argument Antony implies.

\speaker{Antony}
\begin{verse}
Friends, Romans, countrymen, lend me your ears;\\
I come to bury Caesar, not to praise him.\\
The evil that men do lives after them;\\
The good is oft interred with their bones;\\
So let it be with Caesar. The noble Brutus\\
Hath told you Caesar was ambitious:\\
If it were so, it was a grievous fault,\\
And grievously hath Caesar answer'd it.\\
Here, under leave of Brutus and the rest--\\
For Brutus is an honourable man;\\
So are they all, all honourable men--\\
Come I to speak in Caesar's funeral.\\
He was my friend, faithful and just to me:\\
But Brutus says he was ambitious;\\
And Brutus is an honourable man.\\
He hath brought many captives home to Rome\\
Whose ransoms did the general coffers fill:\\
Did this in Caesar seem ambitious?\gloss{Antony offers public benefit as evidence against selfish ambition.}\\
When that the poor have cried, Caesar hath wept:\\
Ambition should be made of sterner stuff:\\
Yet Brutus says he was ambitious;\\
And Brutus is an honourable man.\\
You all did see that on the Lupercal\\
I thrice presented him a kingly crown,\\
Which he did thrice refuse: was this ambition?\gloss{The crown refusal challenges the claim that Caesar simply grasped at kingship.}\\
Yet Brutus says he was ambitious;\\
And, sure, he is an honourable man.\\
I speak not to disprove what Brutus spoke,\\
But here I am to speak what I do know.
\end{verse}

\section*{Independent Analysis}
Answer in complete thoughts. You are testing form first, then the strength of the evidence.

\begin{enumerate}[leftmargin=*]
\item Write Antony's implied conclusion about Caesar's ambition.\studentline
\item Write one possible major premise Antony needs.\studentline\studentline
\item Write one possible minor premise from Antony's evidence.\studentline\studentline
\item Identify S, P, and M in your map.\studentline\studentline
\item Name the mood if you can. Is your conclusion affirmative or negative?\studentline
\item Does Antony's form carry the conclusion, or does it only weaken Brutus's premise? Explain.\studentline\studentline
\item Which speaker gives the cleaner formal argument: Brutus or Antony? Which speaker gives stronger evidence? Explain the difference.\studentline\studentline
\end{enumerate}

\section*{Portfolio Artifact: Julius Caesar Mood-and-Figure Card}
Choose either Brutus's argument or Antony's argument. Complete the card carefully enough that it can be saved in your Logic Portfolio.

\textbf{Speaker and passage:}\studentline
\textbf{Major premise:}\studentline
\textbf{Minor premise:}\studentline
\textbf{Conclusion:}\studentline
\textbf{Mood:}\shortline
\textbf{S, P, and M terms:}\studentline
\textbf{Figure pattern using S, P, and M:}\studentline
\textbf{Form judgment: Does the conclusion follow if the premises are true?}\studentline
\textbf{Fact judgment: Which premise still needs proof?}\studentline\studentline

\section*{Closing Synthesis}
Whately's form test does not make literature smaller. It lets us see more precisely how a speech works. Brutus gives the crowd a strong-looking form: all dangerous ambitious men deserve death; Caesar is such a man; therefore Caesar deserved death. Antony does not merely weep over Caesar. He presses on the middle term, \textit{ambitious}, and asks whether Caesar really belongs under it. The tragedy is not that Rome lacks arguments. The tragedy is that public action can turn on arguments whose form, terms, and premises the crowd has not fully tested.

\end{document}
